\section{产业版图：谁在争夺机器人和物理 AI 的数据闭环}

访谈中提到几股势力：大模型团队、世界模型团队、VLA 团队、本体公司、仿真/合成数据公司。它们不是简单上下游，而是共生关系。谁掌握数据闭环，谁就更可能掌握机器人大脑。

\begin{figure}[H]
\centering
\includegraphics[width=0.95\textwidth]{data-industry-landscape.png}
\caption{数据产业版图：大模型、世界模型、VLA、本体与仿真数据公司。}
\end{figure}

\begin{knowledgebox}{读图：为什么这是共生网络}
大模型团队提供基座和基础设施；世界模型团队理解物理预测；VLA 团队连接动作；本体公司拥有真实机器人；仿真公司提供规模化场景。没有任何一方能单独完成全部数据闭环。
\end{knowledgebox}

\subsection{本体公司 vs 大脑公司}

自动驾驶里，拥有车队的公司天然拥有端侧数据；机器人里，大规模端侧部署尚未出现，因此本体公司未必自动拥有数据优势。大脑公司如果掌握仿真和基础设施，也可能先形成能力；本体公司若能形成真实数据闭环，也会非常关键。

\begin{warningbox}{自动驾驶经验不能直接平移到机器人}
自动驾驶有相对标准化的道路环境和大规模车队；机器人本体、任务和场景更分散。数据引擎需要重新设计，不能照搬 FSD 路线。
\end{warningbox}

\subsection{本章小结}

机器人数据产业会是共生网络。模型、世界模型、VLA、本体和仿真数据公司都可能成为关键节点，关键看谁能形成闭环。

